\section{More Examples of Ring}
\begin{eg}
    $\Z/n\Z$ is an example of a commutative ring because $(\Z/n\Z, +)$ is an abelian group and multiplication modulo $n$ satisfies associativity and commutativity and distributes over addition.
\end{eg}
\begin{eg}
    $\Z/6\Z$ is not a division ring, e.g., $\overline{3}$ doesn't have a multiplicative inverse.
\end{eg}
\begin{eg}[Finite Field]
$\Z/5\Z$ is a division ring. Also it is commutative ring, hence, it is a field.\\
Notice, for all prime numbers $p$, $F_p := \Z/p\Z$ is a field. Number of its elements is finite, hence, it is an example of a finite field.
\end{eg}
\begin{eg}[A ring without multiplicative identity]
    $2\Z$.
\end{eg}
\begin{definition}[Quaternions]
    Let $\mathbb{H}$ to be the set of all quaternions such that 
    \[\mathbb{H} := \{a + bi + cj + dk: a,b,c,d \in R\}.\]
    And define addition $+$ via-
    \[(a_1 + b_1i + c_1j + d_1k)(a_2+b2_i+c_2i + d_2k) := (a_1+a_2) + (b_1 + b_2)i + (c_1+c_2)j + (d_1 + d_2)k.\]
    And multiply using distributive laws subject to the relations:
    \begin{align*}
        &i^2 = j^2 = k^2 = -1\\
        &ij = -ji = k\\
        &jk = -kj = i\\
        &ki = -ik = j
    \end{align*}
\end{definition}
\begin{eg}
The set of all quaternions $\mathbb{H}$ forms a non-commutative division ring.
\end{eg}
\begin{explanation}
\textbf{Quaternions form a ring:}

\begin{itemize}
    \item \textbf{Addition}: The addition of two quaternions is componentwise and is associative and commutative because real numbers are commutative and associative. The additive identity is \( 0 = 0 + 0i + 0j + 0k \), and the additive inverse of \( q = a + bi + cj + dk \) is \( -q = -a - bi - cj - dk \).
    \item \textbf{Multiplication}: The multiplication of quaternions is defined by the distributive law, and is associative.
\end{itemize}
\textbf{Quaternions form a division ring:}
The identity element is \( 1 = 1 + 0i + 0j + 0k \), and every quaternion \( q = a + bi + cj + dk \) has a multiplicative inverse, as shown below. For any non-zero quaternion \( q = a + bi + cj + dk \), the inverse is given by:
\[
q^{-1} = \frac{a - bi - cj - dk}{a^2 + b^2 + c^2 + d^2},
\]
which satisfies \( q q^{-1} = 1 \), as can be verified by multiplying.\\
\textbf{Quaternions form a non-commutative ring:}
By definition, $ij = -ji$, hence, multiplication is not commutative.
\end{explanation}
\begin{eg}[Workout Example]
    $a = 1 + 2j + 3k$ and $b = j + k$.
    \begin{align*}
        a + b &= (1 + 2j + 3k) + (j + k)\\
              &= 1 + (2+1)j + (3+1)k\\
              &= 1 + 3j + 4k.
    \end{align*}
    \begin{align*}
        a \times b &= (1 + 2j + 3k) \times (j + k)\\
                   &= (1 + 2j + 3k)j + (1 + 2j + 3k)k\\
                   &= j + 2j^2 + 3kj + k + 2jk + 3k^2\\
                   &= j + 2(-1) + 3(-i) + k + 2i + 3(-1)\\
                   &= j -2 -3i + k +2i - 3\\
                   &= -5 - i + j + k.
    \end{align*}
\end{eg}
\section{Properties of Ring}
\begin{prop} If $R$ is a ring, then for all $a, b \in R$,
    \begin{enumerate}
        \item $0\cdot a = a \cdot 0 = 0$.
        \item $(-a)b = a(-b) = -(ab)$.
        \item $(-a)(-b) = ab$.
        \item The multiplicative identity is unique.
    \end{enumerate}
\end{prop}
\begin{propproof}
    \begin{enumerate}
        \item \begin{align*}
            (0 + 0) a &= 0\cdot a + 0 \cdot a\\
            \implies 0 \cdot a &= 0\cdot a + 0 \cdot a\\
            \implies 0 &= 0 \cdot a\\
            \implies 0 \cdot a &= 0.
        \end{align*}
        Similarly, $a \cdot 0 = 0$.
        \item \begin{align*}
            (a - a) b &= 0\\
            \implies ab + (-a)b&= 0\\
            \implies (-a)b &= -(ab).
        \end{align*}
        Similarly, $a(-b) = -(ab)$.
        \item Using previous result, $(-a)(-b) = -(a(-b)) = -(-(ab)) = ab.$
        \item Let $1$ and $1'$ be two distinct multiplicative identity. Then for $a \in R$, $a \cdot 1 = a = a \cdot 1'$ which implies $1 = 1'$ (contradiction!). Hence, multiplicative identity has to be unique.
    \end{enumerate}
\end{propproof}
\section{Subring}
\begin{definition}[Subring]
    A subset $S$ of ring $R$ is a \textit{subring} if it is a subgroup of $R$ under addition and if is closed under multiplication.
\end{definition}
\begin{eg}
    $2\Z \subseteq \Z \subseteq \Q \subseteq \R \subseteq \C \subset \mathbb{H}$.
\end{eg}
\section{Characteristic}
\begin{definition}[Characteristic]
    Let $R$ be a ring with identity $1$. If there exists a positive integer $n$ such that $1 + 1 + \cdots + 1 = n1 = 0$, then $n$ is called the characteristic of $R$, and denoted by char($R$). If no such positive integer exists, then $R$ is said to be characteristic zero.
\end{definition}
\begin{eg}
    In $R = \Z/n\Z$, char$(R) = 0$.
\end{eg}